Sparse identification of nonlinear dynamics (SINDy) needs time derivatives that noisy measurements do not provide. We compare five ways of obtaining them, with an identical degree-2 library and a sequentially thresholded least squares (STLSQ) back end, on simulated Lorenz-63 at six noise levels (0--10\% of the signal standard deviation, 20 trajectories per level): finite differences, Savitzky--Golay, a smoothing spline, total-variation (TV) regularised differentiation, and a weak (integral) form that avoids differentiating the data. All five are reimplemented, with smoothing parameters and thresholds tuned per noise level on five separate trajectories. The hypothesis that the weak form recovers the exact support more often than every derivative-based method at every noise level of at least 1\% is refuted: at 1\% it ties three baselines at 1.00, at 2\% (1.00 against 0.80--0.90) and at 5\% (0.65 against 0.00--0.40) its lead is statistically significant against only some baselines, and at 10\% no system exceeds 0.10. The weak form has the lowest mean coefficient error at every level, significantly against all four baselines only at 0\% and 10\%. The hypothesis that smoothing lowers the coefficient error of finite differences is refuted by the registered test: at 5\% all three smoothers have a higher mean error than finite differences. This verdict depends on the finite-difference threshold, which the tuner's tie rule fixed after a first round of results; at equal thresholds only Savitzky--Golay is above finite differences throughout, and TV is below. A threshold grid on the test trajectories shows that up to 1\% noise every method recovers every support at a threshold of 0.4, and that the low-noise ordering of the methods reversed when the tie rule was changed; at 5\% the weak form's rate is above the best threshold of each baseline (0.65 against at most 0.40). This is a one-system, 20-trial CPU study with oracle-assisted tuning.
